%%%%

%%%   Самарский государственный технический университет

%%%   Кафедра прикладной математики и информатики

%%%

%%%   Преамбула для статьи в журнал "Вестник Самарского государственного

%%%   технического университета. Серия: «Физико-математические науки»"

%%%   Ориентирована под MikTEx 2.4 for Win32

%%%

%%%   Хотя сам журнал собирается под GNU/Linux на дистрибутиве TeXLive



\documentclass[11pt,a4paper,twoside]{article}



\usepackage[cp1251]{inputenc}

\usepackage[T2A]{fontenc}

\usepackage[english,russian]{babel}

\usepackage{samgtubib}

\usepackage[dvips]{graphicx}

\usepackage[multiple]{footmisc}



\usepackage{samgtu}

\renewcommand{\VestnikYear}{2009} % Год издания Вестника

\renewcommand{\VestnikNumber}{2\,(19)} % Номер Вестника

\renewcommand{\VestnikMonth}{Ноябрь} % Месяц выхода Вестника

\renewcommand{\VestnikData}{01.11.09} % Дата подписания Вестника в печать

\renewcommand{\VestnikUPL}{26,64} % Усл. печ. л.

\renewcommand{\VestnikUIL}{25,73} % Уч.-изд. л.

\renewcommand{\VestnikRegNumber}{479/09} % Регистрационный номер

\renewcommand{\VestnikOrder}{994} % Номер заказа






%%
%\usepackage{pst-barcode}
%\usepackage{auto-pst-pdf}
%
%\usepackage{calrsfs}
%%%
%%%\newcommand{\totop}[1]{\raisebox{6pt}[1pt][2pt]{#1}} 
%%%\newcommand{\tottop}[1]{\raisebox{12pt}[1pt][2pt]{#1}} 
%%%\newcommand{\totttop}[1]{\raisebox{18pt}[1pt][2pt]{#1}} 
%%%\newcommand{\tottttop}[1]{\raisebox{24pt}[1pt][2pt]{#1}} 
%%%\newcommand{\totttttop}[1]{\raisebox{30pt}[1pt][2pt]{#1}} 
%%%\newcommand{\tobot}[1]{\raisebox{-6pt}[1pt][2pt]{#1}} 
%%%\newcommand{\tobbot}[1]{\raisebox{-12pt}[1pt][2pt]{#1}} 
%%%\newcommand{\tobbbot}[1]{\raisebox{-18pt}[1pt][2pt]{#1}}
%%
%%
%%%\samgtupubl % для генерации pdf, отдаваемого в типографию СамГТУ
%\pdfcrop % обрезка pdf для размещения на math-net.ru
%\draft % На корректуре
%\filllastpage % Последняя страница не заполнена не полностью
%%%\briefcomm % Статья является кратким сообщением

\vsgtu{1387}

\FirstPage{117}
\LastPage{135}
%
%\graphicspath{{./00PIC/}}
%
%\begin{document}

\hfill \begin{pspicture}(1in,1in)
\psbarcode{http://dx.doi.org/10.14498/vsgtu1387}{}{qrcode}
\end{pspicture}
\vspace{-1in}



%\def\I{{\rm I}}
%\def\cl{{C}\!\ell}
\def\R{{\Bbb R}}
\def\C{{\Bbb C}}
\def\F{{\Bbb F}}
\def\N{{\Bbb N}}

\def\Tr{\mathop{\rm Tr}}
\def\cen{\mathop{\rm cen}}
\def\diag{\mathop{\rm diag}}

\def\Adj{\mathop{\rm Adj}}
\def\First{\mathop{\rm First}}
\def\Last{\mathop{\rm Last}}
\def\Even{\mathop{\rm Even}}
\def\Odd{\mathop{\rm Odd}}
\newcommand{\st}{\stackrel}


% Номер УДК
\udk{517.958}%Обработка текста. Подготовка текстов : Языки программирования. Компьютерные языки
\msc{15A66, 70S15}% Коды MSC см. http://www.ams.org/msc/

\rutitle{Свертки по рангам и кватернионным типам в~алгебрах Клиффорда}% Полный заголовок
{Свертки по рангам и кватернионным типам в алгебрах Клиффорда}% Заголовок, который идёт в верхний колонтитул



\ruauthor{Д.~С.~Широков\inst{1,2}}{Ш\,и\,р\,о\,к\,о\,в~Д.~С.}

\ruaffil{\ruippi.\and \rubauman.}

\enaffil{\enippi. \and 
\enbauman.}% Место работы каждого из авторов на англ. языке
%% Если несколько, то так  \enginstitute{ Организация 1 \and Организация 2  \and Организация 3}

%\email{E-mails}{} %E-mail's авторов

\ruauthordetails{1}{\fullauthor{\rupid{52747}{Дмитрий Сергеевич Широков}}\regalia{к.ф.-м.н.; \mem{dm.shirokov@gmail.com}} \position{научный сотрудник} \fdept{лаб. 7 <<Обработка биоэлектрической информации>>\inst{1}}\lposition{доцент} \dept{каф. ФН-1 <<Высшая мате\-ма\-тика>>\inst{2}}
}

%
\enauthordetails{1} {\fullauthor{\enpid{52747}{Dmitry S. Shirokov}} \regalia{Cand. Phys. \& Math. Sci.; \mem{dm.shirokov@gmail.com}} \position{Scientific Researcher} 
\fdept{Lab. 7 ``Bioelectric Information Processing''\inst{1}}
\lposition{Assistant Professor} \dept{Dept. of Higher Mathematics\inst{2}}}

%\otherinf{\fullauthor{}\regalia{к.ф.-м.н.}\position{научный сотрудник} \dept{лаб. 7 ИППИ РАН}\position{доцент}
%\dept{каф. ФН-1 '' Высшая математика'' МГТУ им. Н.Э.Баумана}
%}

\rukeywords{алгебры Клиффорда, свертки, операции проецирования, кватернионный тип}

\ruabstract{В работе рассмотрены выражения в вещественных и комплексных алгебрах Клиффорда, называемые свертками или усреднениями. Свертка берется от произвольного элемента алгебры Клиффорда, при этом ведется суммирование по различным элементам фиксированного базиса алгебры Клиффорда. Рассмотрены четные и нечетные свертки, свертки по рангам и свертки по кватернионным типам. Представлена связь сверток с операциями проецирования на выделенные подпространства алгебры Клиффорда "--- четное и нечетное подпространство, подпространства фиксированных рангов и подпространства фиксированных кватернионных типов. С помощью метода сверток дано решение различных систем коммутаторных уравнений в алгебрах Клиффорда. Особое внимание уделено двум частным случаям "--- случаям коммутатора и антикоммутатора. Полученные результаты могут применяться при изучении различных уравнений теории поля "--- уравнений Янга"--~Миллса, простейшего полевого уравнения и других.}
%
\entitle{Contractions on ranks and quaternion types in~Clifford algebras}

\enauthor{D.~S.~Shirokov\inst{1,2}}{S\,h\,i\,r\,o\,k\,o\,v~D.~S.}
%
%


%
\enabstract{In this paper we consider expressions in real and complex Clifford algebras, which we call contractions or averaging. We consider contractions of arbitrary Clifford algebra element. Each contraction is a sum of several summands with different basis elements of Clifford algebra. We consider even and odd contractions, contractions on ranks and contractions on quaternion types. We present relation between these contractions and projection operations onto fixed subspaces of Clifford algebras --- even and odd subspaces, subspaces of fixed ranks and subspaces of fixed quaternion types. Using method of contractions we present solutions of system of commutator equations in Clifford algebras. The cases of commutator and anticommutator are the most important. These results can be used in the study of different field theory equations, for example, Yang--Mills equations, primitive field equation and others.
}

\enkeywords{Clifford algebras, contractions, projection operations, quaternion type}

\date{15/XII/2014} % Дата поступления статьи в редакцию.
\revisiondate{17/II/2015} % В окончательном виде
\accepteddate{25/II/2015}


%%%%%%%%%%%%%%%%%%%%%%%%%%%%%%%%%%%%%%%%%%%%%%%%%%%%%%%%%%%%%%%%%%%%%%%%%%%%%%%%%%%%


\makerutitle

\Section[n]{Введение}
Алгебра Клиффорда была предложена в 1878 году У.~Клиффордом \cite{shir:Clifford}. 
В~своих исследованиях он объединил идеи, связанные с кватернионами Гамильтона~\cite{shir:Hamilton} и~внешней алгеброй Грассмана~\cite{shir:Grassmann}. 
В~дальнейшем алгебра Клиффорда развивалась усилиями многих известных математиков "--- Р.~Липшицем~\cite{shir:Lipsch}, Э.~Картаном, Э.~Уиттом, К.~Шевалле~\cite{shir:Chevall}, М.~Риссом и~другими. 
Существенное влияние на развитие теории алгебр Клиффорда оказало открытие уравнения Дирака для электрона в~1928~году~\cite{shir:Dirac}. 
В~настоящее время алгебры Клиффорда широко применяются в~различных разделах современной математики и физики "--- теории поля~\cite{shir:Hestenes, shir:Marchuk}, робототехнике, небесной механике, обработке сигналов и~изображений, вычислительной технике, химии, геометрии.

В настоящей статье мы рассматриваем выражения в алгебре Клиффорда вида
$$
\sum_{A\in S}e_A U e^A,\qquad e_A=(e^A)^{-1},
$$
где $e^A$ "--- элементы базиса алгебры Клиффорда, $S$ есть подмножество множества $\rm I$ всех упорядоченных мультииндексов $A$ длины от $0$ до $n$. 
Будем называть такие выражения свертками, или усреднениями, в алгебре Клиффорда.

Отметим, что метод сверток напрямую связан с методом усреднения в~теории представлений конечных групп \cite{shir:tr1, shir:tr2}.

В работе автора \cite{shir:aacacontr} изучены полные свертки (случай $S=\rm I$), простые свертки (множество $S$ состоит из одного элемента) и свертки по сопряженным наборам мультииндексов.

В~настоящей работе продолжено изучение сверток в алгебрах Клиффорда. 
Рассматриваются четные и нечетные свертки (когда множество $S$ содержит мультииндексы четной или нечетной длины), свертки по рангам (участвуют мультииндексы фиксированной длины), свертки по кватернионным типам. 
Доказываются теоремы о~связи различных сверток с~проекциями на выделенные подпространства алгебры Клиффорда.

Даны решения систем коммутаторных уравнений в алгебрах Клиффорда
$$
e^A X+\epsilon X e^A=q^A,\qquad A\in S\subseteq {\rm I},\qquad \epsilon\in\C\setminus\{0\}
$$
для неизвестного элемента $X\in{{C}\!\ell}(p,q)$ и известных элементов $q^A\in{{C}\!\ell}(p,q)$.



Техника сверток напрямую связана с изучением различных уравнений теории поля. 
Так, в работе \cite{shir:YM} рассматривается простейшее полевое уравнение и с помощью техники сверток найдено его общее решение. 
Там же с~помощью техники сверток предложен новый класс решений уравнений Янга"--~Миллса.

В работе \cite{shir:gencontr} рассмотрены обобщенные свертки, построенные по двум наборам антикоммутирующих элементов алгебры Клиффорда. 
С~помощью обобщенных сверток дано обобщение теоремы Паули~\cite{shir:Pauli} на случай алгебры Клиффорда~\cite{shir:gencontr} и~решен ряд вопросов о~связи спинорных и~ортогональных групп [\citen{shir:Pauli1}--\citen{shir:Pauli3}].

\smallskip

\Section{Вещественные и комплексные алгебры Клиффорда, кватернионный тип элемента}
Рассмотрим комплексную алгебру Клиффорда ${{C}\!\ell}(p,q)$ (или вещественную ${{C}\!\ell}^\R(p,q)$), где $p+q=n$, $n\geq1$.
Построение алгебры Клиффорда подробно приведено в \cite{shir:Marchuk:Shirokov} и \cite{shir:Lounesto}. 
Будем называть размерностью алгебры Клиффорда ${{C}\!\ell}(p,q)$ число $n$, хотя ее размерность как линейного пространства равна $2^n$.

Единичный элемент обозначим через $e$, а генераторы алгебры Клиффорда ${{C}\!\ell}(p,q)$ через $e^a$, $a=1,\ldots,n$. Генераторы удовлетворяют определяющим соотношениям
$$
e^a e^b+e^b e^a=2\eta^{ab}e,
$$
где $\eta=\|\eta^{ab}\|=\|\eta_{ab}\|=\diag(1, \ldots, 1, -1, \ldots, -1)$ "--- диагональная матрица с $p$ единицами и $q$ минус единицами на диагонали. 
Элементы
$$
e^{a_1\ldots a_k}=e^{a_1}\ldots e^{a_k},\qquad a_1<\ldots<a_k,\qquad k=1,\ldots,n,
$$
вместе с единичным элементом $e$ образуют базис алгебры Клиффорда. Всего имеется $2^n$ элементов базиса.

Через $\rm I$ будем обозначать множество мультииндексов длины от $0$ до $n$
\begin{equation}
{\rm I}=\{-,\, 1,\, \ldots,\, n,\, 12,\, 13,\, \ldots,\, 1\ldots n\},\label{shirokov:Inab}
\end{equation}
где через $-$ обозначен пустой мультииндекс, соответствующий единичному элементу алгебры Клиффорда. 
Итак, мы имеем базис алгебры Клиффорда $\{e^A,\, A\in{\rm I} \}$, где $A$ есть произвольный упорядоченный мультииндекс. Обозначим длину мультииндекса $A$ через $|A|$.

Будем рассматривать различные подмножества $S\subseteq \rm I$:
$${\rm I}_ {\Even}=\{A\in{\rm I},\, |A| - \mbox{четно}\},\qquad {\rm I}_ {\Odd}=\{A\in{\rm I},\, |A| - \mbox{нечетно}\},$$
$${\rm I}_ {k}=\{A\in{\rm I},\quad |A|=k\},\qquad k=0, 1, \ldots, n,$$
$${\rm I}_ {\overline k}=\{A\in{\rm I},\quad |A|=k\mod 4\},\qquad k=0, 1, 2, 3.$$

Индексы опускаются и поднимаются с помощью матрицы $\eta$, т.е. $e_a=\eta_{ab}e^{b}$, $e^a=\eta^{ab}e_b$. 
Мы пользуемся соглашением Эйнштейна о~суммировании по повторяющемуся нижнему и верхнему индексу. 
Имеем
$$e_{a_1\ldots a_k}=\eta_{a_1 b_1}\ldots \eta_{a_k b_k} e^{b_k}\ldots e^{b_1}=e_{a_k}\ldots e_{a_1}=(e^{a_1\ldots a_k})^{-1},\quad a_1<\ldots< a_k.$$


Произвольный элемент алгебры Клиффорда $U\in{{C}\!\ell}(p,q)$ может быть записан в виде
\begin{equation}
U=ue+u_a e^a+\sum_{a_1<a_2}u_{a_1 a_2}e^{a_1 a_2}+\ldots+u_{1\ldots n}e^{1\ldots n}=u_A e^A,\label{shirokov:U}
\end{equation}
где $\{u_A\}=\{u, u_a, u_{a_1 a_2}, \ldots, u_{1\ldots n}\}$  "--- произвольные комплексные (или вещественные) константы.

Обозначим через ${{C}\!\ell}_k(p,q)$ векторные подпространства, натянутые на элементы
$e^{a_1\ldots a_k}$. Элементы ${{C}\!\ell}_k(p,q)$ называются элементами {\it ранга $k$}. Имеем
\begin{equation}
{{C}\!\ell}(p,q)=\bigoplus_{k=0}^{n}{{C}\!\ell}_k(p,q),\qquad \dim{{C}\!\ell}_k(p,q)=C_n^k=\frac{n!}{k!(n-k)!}.\label{shirokov:ranks}
\end{equation}
Рассмотрим проекционные операторы на подпространства ${{C}\!\ell}_k(p,q)$
\begin{equation}
\pi_k: {{C}\!\ell}(p,q)\to {{C}\!\ell}_k(p,q),\qquad \pi_k(U)=\sum_{a_1<\ldots<a_k}u_{a_1\ldots a_k}e^{a_1 \ldots a_k}.\label{shirokov:proj}
\end{equation}

Центром алгебры Клиффорда 
$${\rm cen}{{C}\!\ell}(p,q)=\{U\in {{C}\!\ell}(p,q) \,|\, UV=VU\quad \forall V\in{{C}\!\ell}(p,q)\}$$  является следующее множество:
\begin{equation}
\cen{{C}\!\ell}(p,q)=\left\lbrace\begin{array}{ll}
{{C}\!\ell}_0(p,q), & \mbox{если $n$ четно};\\
{{C}\!\ell}_0(p,q)\oplus {{C}\!\ell}_n(p,q), & \mbox{если $n$ нечетно.}
\end{array}
\right.
\end{equation}

Алгебра Клиффорда ${{C}\!\ell}(p,q)$ является супералгеброй. 
Она может быть представлена в виде прямой суммы четного и нечетного подпространства:
$${{C}\!\ell}(p,q)={{C}\!\ell}_{\Even}(p,q)\oplus{{C}\!\ell}_{\Odd}(p,q),\quad \dim{{C}\!\ell}_{\Even}(p,q)=\dim{{C}\!\ell}_{\Odd}(p,q)=2^{n-1},$$
$${{C}\!\ell}_{\Even}(p,q)=\bigoplus_{k - \text{even}}{{C}\!\ell}_k(p,q),\qquad {{C}\!\ell}_{\Odd}(p,q)=\bigoplus_{k - \text{odd}}{{C}\!\ell}_k(p,q).$$

Представим алгебру Клиффорда ${{C}\!\ell}^\R(p,q)$ как прямую сумму четырех подпространств {\it кватернионных типов} 0, 1, 2 и 3 (см. [\citen{shir:quat1}--\citen{shir:quat3}]):
\begin{equation}
{{C}\!\ell}^\R(p,q)={{C}\!\ell}^\R_{\overline 0}(p,q)\oplus{{C}\!\ell}^\R_{\overline 1}(p,q)\oplus
{{C}\!\ell}^\R_{\overline 2}(p,q)\oplus{{C}\!\ell}^\R_{\overline 3}(p,q),\label{shirokov:kv}
\end{equation}
где
\begin{equation}
{{C}\!\ell}^\R_{\overline m}(p,q)=\bigoplus_{k=m\mod 4}  {{C}\!\ell}^\R_k(p,q),\qquad m=0, 1, 2, 3.
\label{shirokov:kv-1}
\end{equation}

Соответствующие операции проецирования на подпространства алгебры Клиффорда фиксированного кватернионного типа обозначены через 
$$\pi_{\overline k}:{{C}\!\ell}(p,q)\to {{C}\!\ell}_{\overline k}(p,q),\qquad k=0, 1, 2, 3.$$

Обозначая для краткости ${{C}\!\ell}^\R_{\overline k}(p,q)$ через $\overline{\textbf{k}}$, можем записать  свойства для этих подпространств относительно действия коммутатора $$[U,V]=UV-VU$$ и антикоммутатора $$\{U,V\}=UV+VU$$ (см. более подробно \cite{shir:quat1,shir:quat2}):
\begin{eqnarray*}
&&[\overline{\textbf{k}},\overline{\textbf{k}}]\subseteq\overline{\textbf{2}},\quad \{\overline{\textbf{k}},\overline{\textbf{k}}\}\subseteq\overline{\textbf{0}},\quad [\overline{\textbf{k}},\overline{\textbf{2}}]\subseteq\overline{\textbf{k}}, \quad \{\overline{\textbf{k}},\overline{\textbf{0}}\}\subseteq\overline{\textbf{k}},\qquad k=0, 1, 2, 3 \nonumber;\label{shirokov:1}\\
&&[\overline{\textbf{0}},\overline{\textbf{1}}]\subseteq\overline{\textbf{3}}, \quad  [\overline{\textbf{0}},\overline{\textbf{3}}]\subseteq\overline{\textbf{1}}, \quad [\overline{\textbf{1}},\overline{\textbf{3}}]\subseteq\overline{\textbf{0}},\quad \{\overline{\textbf{1}},\overline{\textbf{2}}\}\subseteq\overline{\textbf{3}},  \quad \{\overline{\textbf{1}},\overline{\textbf{3}}\}\subseteq\overline{\textbf{2}}, \quad \{\overline{\textbf{2}},\overline{\textbf{3}}\}\subseteq\overline{\textbf{1}}\nonumber.
\end{eqnarray*}
Комплексную алгебру Клиффорда ${{C}\!\ell}(p,q)$ можно представить в виде прямой суммы восьми следующих подпространств:
$${{C}\!\ell}(p,q)=\overline{\textbf{0}}\oplus\overline{\textbf{1}}\oplus\overline{\textbf{2}}\oplus\overline{\textbf{3}} \oplus i\overline{\textbf{0}}\oplus i\overline{\textbf{1}}\oplus i\overline{\textbf{2}}\oplus i\overline{\textbf{3}}.$$

\smallskip

\hypertarget{shir:par2}{}
\Section{Свертки по рангам}
Рассмотрим следующую {\it свертку генераторов}\cite{shir:Marchuk:Shirokov} в алгебре Клиффорда ${{C}\!\ell}(p,q)$:
$$F_1(U)=\sum_{A\in {\rm I}_ 1} e_A U e^A=e_a U e^a,\qquad U\in{{C}\!\ell}(p,q).$$

\smallskip

%\label{shirokov:theoremSvGen}
\hypertarget{shirokov:theoremSvGen}{}
\begin{theorem}[1~\cite{shir:Marchuk:Shirokov}]Для произвольного элемента $U\in{{C}\!\ell}(p,q)$ имеем
$$F_1(U)=e_a U e^a=\sum_{k=0}^n (-1)^k (n-2k)\pi_k(U).$$
\end{theorem}

\smallskip

Отметим, что в работе \cite{shir:YM} мы представили связь между сверткой генераторов и проекциями на подпространства алгебры Клиффорда фиксированного ранга. 
Используя эту связь, удалось предъявить класс новых решений уравнений Янга"--~Миллса.

%Сформулируем эту теорему.
%Мы используем обозначения $F_1^0(U)=U$, $F_1^1(U)=F(U)$, $F_1^2(U)=F(F(U))$ и т.д.
%Оператор $F_1^l: {{C}\!\ell}(p,q)\to{{C}\!\ell}(p,q)$ называется {\it сверткой генераторов порядка $l$}.
%
%\begin{theorem} \cite{YM} Рассмотрим произвольный элемент алгебры Клиффорда $U\in{{C}\!\ell}(p,q)$, $n=p+q$. Имеем
%\begin{eqnarray}
%\pi_k(U)&=&\sum_{l=0}^n b_{kl} F^l(U),\qquad \mbox{если $n$ четно,}\label{shirokov:svpr1}\\
%\pi_{k, n-k}(U)&=&\sum_{l=0}^{\frac{n-1}{2}}g_{kl} F^l(U),\qquad \mbox{если $n$ нечетно,\label{shirokov:svpr2}}
%\end{eqnarray}
%где $B=||k_{kl}||$ - обратная матрица к матрице $A_{(n+1) \times (n+1)}=||a_{kl}||$, $a_{kl}=(\lambda_{l-1})^{k-1}$, $G=||g_{kl}||$ - обратная матрица к матрице $D_{\frac{n+1}{2} \times \frac{n+1}{2}}=||d_{kl}||$, $d_{kl}=(\lambda_{l-1})^{k-1}$ и $\lambda_k=(-1)^k(n-2k)$.
%\end{theorem}

Теперь рассмотрим более общие выражения (совпадающие с предыдущими при $m=1$) "--- свертки по фиксированным рангам
$$F_m(U)=\sum_{A\in{\rm I}_ m}e_A U e^A,\qquad {\rm I}_ {m}=\{A\in{\rm I},\quad |A|=m\}\qquad m=0, 1, \ldots ,n.$$

\smallskip
%\label{shirokov:theoremSvBas} 
\hypertarget{shirokov:theoremSvBas}{}
\begin{theorem}[2]Для произвольного элемента $U\in{{C}\!\ell}(p,q)$ имеем
\begin{gather}
F_m(U)=\sum_{A\in{\rm I}_ m}e_A U e^A=\sum_{k=0}^n(-1)^{km}\biggl(\sum_{i=0}^{m}(-1)^iC_k^i C_{n-k}^{m-i}\biggr)\pi_k(U).\label{shirokov:svert}
\end{gather}
\end{theorem}

\smallskip

Заметим, что, используя свертку Вандермонда $C_n^m=\sum_{i=0}^{m}C_k^i C_{n-k}^{m-i}$, можем записать
$$\sum_{i=0}^{m}(-1)^iC_k^i C_{n-k}^{m-i}=C_n^m-2\sum_{i - \text{odd}}C_k^i C_{n-k}^{m-i}=-C_n^m+2\sum_{i - \text{even}}C_k^i C_{n-k}^{m-i}.$$

\smallskip

\begin{proof} В~силу линейности достаточно доказать следующую формулу для произвольного элемента базиса алгебры Клиффорда:
$$\sum_{b_1<\ldots<b_m}(e^{b_1 \ldots b_m})^{-1}e^{a_1\ldots a_k} e^{b_1\ldots b_m}=
(-1)^{km}\biggl (\sum_{i=0}^{m}(-1)^iC_k^i C_{n-k}^{m-i}\biggr)e^{a_1\ldots a_k}.$$
Пусть $i$ "--- число совпадающих индексов в упорядоченных наборах $a_1\ldots a_k$ и $b_1\ldots b_m$. 
Пусть для определенности $k\geq m$. 
Для каждого $i$ будет $C_k^i C_{n-k}^{m-i}$ различных наборов $b_1\ldots b_m$. 
В~каждом случае при перестановке местами элементов $e^{b_1 \ldots b_m}$ и $e^{a_1\ldots a_k}$ будем получать коэффициент $(-1)^{km-i}$. 
Из этих соображений получаем нужную формулу. 
При $k<m$ получаем ту же формулу, т.к. $C_k^i=0$ для $i>k.$
\end{proof}

\smallskip

Приведем несколько примеров ($m=2, 3, 4, n$):
\begin{eqnarray*}
e_{b_1 b_2}U e^{b_1 b_2}&=&\sum_k(C^2_n-2k(n-k))\pi_k(U),\nonumber\\
e_{b_1 b_2 b_3}U e^{b_1 b_2 b_3}&=&\sum_k(-1)^k(C_n^3-2(k C^2_{n-k}+C_k^3))\pi_k(U),\nonumber\\
e_{b_1 b_2 b_3 b_4}U e^{b_1 b_2 b_3 b_4}&=&\sum_k(C_n^4-2(k C^3_{n-k}+(n-k)C_k^3))\pi_k(U),\nonumber\\
e_{1\ldots n}Ue^{1\ldots n}&=&\sum_k(-1)^{k(n+1)}\pi_k(U).\nonumber
\end{eqnarray*}

Рассмотрим систему коммутаторных уравнений по рангам, сначала в частном случае "--- по генераторам. 
Верна следующая теорема.

\smallskip

\begin{theorem}[3]Пусть элемент $X\in {{C}\!\ell}(p,q)$ удовлетворяет следующей системе $n$ уравнений для некоторых элементов $q^{a}\in{{C}\!\ell}(p,q){:}$
\begin{equation}
[e^{a}, X]=q^{a},\qquad a=1, 2,\ldots, n.
\end{equation}
Тогда система либо не имеет решения$,$ либо имеет единственное с точностью до элемента центра решение вида
\begin{equation}
X=
\left\lbrace
\begin{array}{ll}
\displaystyle
\sum_{k=1}^{n}\frac{\pi_k( q^{a} e_{a})}{(-1)^{k}(n-2k)-n}+ U_0,  & \mbox{если $n$ четно},\\[2mm]
\displaystyle
\sum_{k=1}^{n-1}\frac{\pi_k( q^{a} e_{a})}{(-1)^{k}(n-2k)-n}+ U_0 +U_n,  & \mbox{если $n$ нечетно},
\end{array}
\right.
\end{equation}
где $U_0,$ $U_n$ "--- произвольные элементы рангов $0$ и $n$ соответственно\rm.
\end{theorem}

\smallskip

\begin{proof} Домножим уравнение справа на $e_a$ и просуммируем: 
$$e^{a} X e_{a}-X e^{a} e_{a}=q^{a} e_{a}.$$
Отсюда получаем
$$\sum_{k=0}^{n}(-1)^{k}(n-2k)\pi_k(X)- n X=q^{a} e_{a}.$$
Расписывая $X=\sum_{k=0}^n \pi_k(X)$, получаем
$$\sum_{k=0}^{n}((-1)^{k}(n-2k)-n)\pi_k(X)=\sum_{k=0}^n \pi_k(q^{a} e_{a}).$$
Выражение $(-1)^{k}(n-2k)-n$ равняется нулю только при $k=0$ и при $k=n$, если $n$ нечетно. 
Отсюда получаем утверждение теоремы.
\end{proof}

\smallskip

Теперь рассмотрим более общий случай, а именно систему $C_n^m$ уравнений для неизвестного $X\in{{C}\!\ell}(p,q)$ и известных $q^{a_1 \ldots a_m}\in{{C}\!\ell}(p,q)$ (число $m$ фиксировано) вида
\begin{equation}
e^{a_1 \ldots a_m} X+\epsilon X e^{a_1 \ldots a_m}=q^{a_1 \ldots a_m},\qquad A={a_1 \ldots a_m}\in {\rm I}_ m,\qquad \epsilon\in\C\setminus\{0\}.
\end{equation}
Утверждение для общего случая довольно громоздко, поэтому опишем лишь сам метод решения таких систем уравнений, который лучше применять уже для конкретных $m$ и $\epsilon$. 
В~частности, при $\epsilon=-1$ и $m=1$ получаем утверждение из предыдущей теоремы.

Домножая справа каждое уравнение системы на соответствующий обратный элемент и суммируя уравнения, получаем
$$e^{a_1 \ldots a_m} X e_{a_1 \ldots a_m}+\epsilon X e^{a_1 \ldots a_m} e_{a_1 \ldots a_m}=q^{a_1 \ldots a_m} e_{a_1 \ldots a_m}.$$
Тогда
$$\sum_{k=0}^{n}(-1)^{km}\sum_{i=0}^m (-1)^i C_k^i C_{n-k}^{m-i}\pi_k(X)+\epsilon X C_n^m=q^{a_1 \ldots a_m} e_{a_1 \ldots a_m}.$$
Подставляя $X=\sum_{k=0}^n\pi_k(X)$, получаем
$$\sum_{k=0}^{n}\biggl((-1)^{km}\sum_{i=0}^m (-1)^i C_k^i C_{n-k}^{m-i}+\epsilon C_n^m\biggr)\pi_k(X)=\sum_{k=0}^n \pi_k(q^{a_1 \ldots a_m} e_{a_1 \ldots a_m}).$$
Далее действуем так же, как при доказательстве предыдущей теоремы. Система либо не будет  иметь решения (в~зависимости от элементов $q^{a_1 \ldots a_m}$), либо будет иметь решение, расписанное через сумму различных проекций. При этом решение будет единственным с точностью до прибавления произвольных элементов некоторых рангов. Эти ранги $k$ определяются тем, при каких $k$ выполнено 
$$(-1)^{km}\sum_{i=0}^m (-1)^i C_k^i C_{n-k}^{m-i}+\epsilon C_n^m=0$$ и зависят, таким образом, от $n$, $m$ и $\epsilon$.

\smallskip

\Section{Четные и нечетные свертки}
В работе \cite{shir:aacacontr} были рассмотрены {\it полные свертки} $F(U)=\frac{1}{2^n} e_A U e^A.$
Они проецируют произвольный элемент алгебры Клиффорда на центр:
\begin{equation}
F(U)=\frac{1}{2^n}e_A U e^A=\left\lbrace
\begin{array}{ll}
\pi_0(U), & \parbox{.2\linewidth}{если $n$ четно;} \\
\pi_0(U)+\pi_n(U), & \parbox{.2\linewidth}{если $n$ нечетно.}
\end{array}
\right.\label{shirokov:poln}
\end{equation}



Теперь рассмотрим  {\it четные} и {\it нечетные свертки}:
$$F_{\Even}(U)=\frac{1}{2^{n-1}}\sum_{A\in{\rm I}_ {\Even}}e_A U e^A,\qquad F_{\Odd}(U)=\frac{1}{2^{n-1}}\sum_{A\in{\rm I}_ {\Odd}}e_A U e^A.$$

\smallskip

%\label{shirokov:theoremSvEvOdd} 
\hypertarget{shirokov:theoremSvEvOdd}{} 
\begin{theorem}[4]Для произвольного элемента алгебры Клиффорда $U\in{{C}\!\ell}(p,q),$ $n=p+q$ имеем
\begin{gather}
F_{\Even}(U)=\frac{1}{2^{n-1}}\sum_{A\in{\rm I}_ {\Even}}e_A U e^A=\pi_0(U)+\pi_n(U),\label{shirokov:evod}\\
F_{\Odd}(U)=\frac{1}{2^{n-1}}\sum_{A\in{\rm I}_ {\Odd}}e_A U e^A=\pi_0(U)+ (-1)^{n+1}\pi_n(U).\nonumber
\end{gather}
Рассматриваемые операторы являются проекторами
$F_{\Even}^2=F_{\Even},$ $F_{\Odd}^2=F_{\Odd}.$ В случае нечетного $n$ имеем $F=F_{\Even}=F_{\Odd}.$
В случае четного $n$ имеем $F=\frac{1}{2}(F_{\Even} + F_{\Odd}).$
\end{theorem}

\smallskip

\begin{proof} В силу линейности достаточно проверить, как действуют рассматриваемые свертки на произвольный элемент алгебры Клиффорда фиксированного ранга
$U_k\in{{C}\!\ell}_k(p,q)$, $k=0, \ldots, n$.
Случай $k=0$ тривиален, т.к. 
$$\sum_{A\in{\rm I}_ {\Even}}e_A e^A=\sum_{A\in{\rm I}_ {\Odd}}e_A e^A=2^{n-1}.$$ 
В случае нечетного $k=n$ элементы ранга $n$ лежат в центре, поэтому $$\sum_{A\in{\rm I}_ {\Even}}e_A e^{1\ldots n}e^A=\sum_{A\in{\rm I}_ {\Odd}}e_A e^{1\ldots n}e^A=2^{n-1}e^{1\ldots n}.$$ 

\pagebreak

\noindent
В случае четного $k=n$ элемент $e^{1\ldots n}$ антикоммутирует со всеми нечетными элементами алгебры Клиффорда и коммутирует со всеми четными элементами алгебры Клиффорда, поэтому
$$\sum_{A\in{\rm I}_ {\Even}}e_A e^{1\ldots n}e^A=-\sum_{A\in{\rm I}_ {\Odd}}e_A e^{1\ldots n}e^A=2^{n-1}e^{1\ldots n}.$$



Для всех остальных рангов $k=1,\ldots, n-1$ можем воспользоваться \eqref{shirokov:svert}, просуммировать по всем четным (или нечетным) $m$, перегруппировать слагаемые и получить
\begin{multline}
\sum_{A\in{\rm I}_ {\Even}}e_A U_k e^A=
\biggl(
\biggl(\sum_{i - \text{even}}C_k^i\biggr)\biggl(\sum_{i - \text{even}}C_{n-k}^i\biggr)- \\ -
\biggl(\sum_{i - \text{odd}}C_k^i\biggr)\biggl(\sum_{i - \text{odd}}C_{n-k}^i\biggr)\biggr)U_k=\\
=(2^{k-1} 2^{n-k-1}-2^{k-1} 2^{n-k-1})U_k=0,
\end{multline}
\begin{multline}
\sum_{A\in{\rm I}_ {\Odd}}e_A U_k e^A=(-1)^k
\biggl(\biggl(\sum_{i - \text{even}}C_k^i\biggr)\biggl(\sum_{i - \text{odd}}C_{n-k}^i\biggr)- \\ - \biggl(\sum_{i - \text{odd}}C_k^i\biggr)\biggl(\sum_{i - \text{even}}C_{n-k}^i\biggr)\biggr)U_k=
\\=(-1)^k (2^{k-1} 2^{n-k-1}-2^{k-1} 2^{n-k-1})U_k=0.
\end{multline}
Теорема доказана.
\end{proof}

\smallskip

\begin{theorem}[5]Для произвольного элемента $U\in{{C}\!\ell}(p,q)$ в случае четного $n=p+q$ имеем
\begin{eqnarray*}
\pi_0(U)&=&\frac{1}{2^n}\biggl(\sum_{A\in{\rm I}_ {\Even}}e_A U e^A+\sum_{A\in{\rm I}_ {\Odd}}e_A U e^A\biggr)=\frac{1}{2^n}e_A U e^A,\\
\pi_n(U)&=&\frac{1}{2^n}\biggl(\sum_{A\in{\rm I}_ {\Even}}e_A U e^A-\sum_{A\in{\rm I}_ {\Odd}}e_A U e^A\biggr).
\end{eqnarray*}
\end{theorem}

\smallskip
\begin{proof} Используя \eqref{shirokov:evod}, получаем утверждение теоремы. \end{proof}

\smallskip

Итак, в случае четного $n$ проекцию элемента на подпространства рангов~$0$ и~$n$ можно получить как линейную комбинацию двух рассматриваемых сверток.
Если $n$ нечетно, то можем получить только проекцию на центр (с помощью любой из трех рассматриваемых сверток):
$$\pi_{center}(U)=\frac{1}{2^{n-1}}\sum_{A\in{\rm I}_ {\Even}}e_A U e^A=\frac{1}{2^{n-1}}\sum_{A\in{\rm I}_ {\Odd}}e_A U e^A=\frac{1}{2^n}e_A U e^A,$$
где $\pi_{center}$ "--- проекция на центр алгебры Клиффорда.

Теперь рассмотрим системы коммутаторных уравнений по четным и нечетным мультииндексам.

\smallskip

\begin{theorem}[6]Пусть элемент $X\in {{C}\!\ell}(p,q)$ удовлетворяет следующей системе $2^{n-1}$ уравнений для некоторых элементов $q^A\in{{C}\!\ell}(p,q){:}$
\begin{equation}
e^A X+\epsilon X e^A=q^A,\quad |A|- \text{even},\qquad \epsilon\in\C\setminus\{0\}.
\end{equation}
Тогда в случае $\epsilon=-1$  $($случай коммутатора$)$
система либо не имеет ре\-ше\-ния$,$ либо имеет решение $($единственное с точностью до прибавления произвольных элементов рангов $0$ и $n){:}$
\begin{equation}
X=-\frac{1}{2^{n-1}}\sum_{|A|-\text{even}}q^A e_A+U_0+U_n.
\end{equation}

В случае $\epsilon\neq -1$ система либо не имеет ре\-ше\-ния$,$ либо имеет единственное решение{\rm:}
\begin{equation}
X=\frac{1}{2^{n-1}\epsilon}\biggl(\sum_{|A|-\text{even}}q^A e_A-\frac{1}{(\epsilon+1)}(\pi_0(q^Ae_A)+\pi_n(q^Ae_A)\biggr).
\end{equation}
\end{theorem}
\smallskip

\begin{proof} Имеем
$$\sum_{|A|-\text{even}}e^A X e_A+\epsilon X \sum_{|A|-\text{even}}e^A e_A=\sum_{|A|-\text{even}}q^A e_A,$$
откуда
$$2^{n-1}(\pi_0(X)+\pi_n(X))+\epsilon 2^{n-1} X =\sum_{|A|-\text{even}}q^A e_A.$$
Расписывая элемент по рангам $X=\sum_{k=0}^n \pi_k(X)$, получаем решение.
\end{proof}


\smallskip
\begin{theorem}[7]Пусть элемент $X\in {{C}\!\ell}(p,q)$ удовлетворяет следующей системе $2^{n-1}$ уравнений для некоторых элементов $q^A\in{{C}\!\ell}(p,q){:}$
\begin{equation}
e^A X+\epsilon X e^A=q^A,\quad |A|- \text{odd},\qquad \epsilon\in\C\setminus\{0\}.
\end{equation}
Тогда в случае $\epsilon=-1$ $($случай коммутатора$)$ система либо не имеет ре\-ше\-ние$,$ либо имеет решение $($единственное с точностью до прибавления элементов заданных ран\-гов$)$ вида
$$X=-\frac{1}{2^{n-1}}\biggl(\sum_{|A|-\text{odd}}q^A e_A-\frac{1}{2}\pi_n(\sum_{|A|-\text{odd}}q^A e_A)\biggr)+U_0$$ в случае четного $n$ и
$$X=-\frac{1}{2^{n-1}}\sum_{|A|-\text{odd}}q^A e_A+U_0+U_n$$ в случае нечетного $n.$

В случае $\epsilon=1$ $($случай антикоммутатора$)$ система либо не имеет ре\-ше\-ния$,$ либо имеет решение вида $($единственное с точностью до прибавления элементов заданных рангов$)$ вида
$$X=\frac{1}{2^{n-1}}\biggl(\sum_{|A|-\text{odd}}q^A e_A-\frac{1}{2}\pi_n\biggl(\sum_{|A|-\text{odd}}q^A e_A\biggr)\biggr)+U_n$$
в случае четного $n$ и
$$X=\frac{1}{2^{n-1}}\biggl(\sum_{|A|-\text{odd}}q^A e_A-\frac{1}{2}\pi_0\biggl(\sum_{|A|-\text{odd}}q^A e_A\biggr)-\frac{1}{2}\pi_n\biggl(\sum_{|A|-\text{odd}}q^A e_A\biggr)\biggr)$$
в случае нечетного $n.$

В случае $\epsilon\neq\pm 1$ система либо не имеет ре\-ше\-ния$,$ либо имеет единственное решение вида
$$X=\frac{1}{2^{n-1}\epsilon}\biggl(\sum_{|A|-\text{odd}}q^A e_A-\frac{1}{\epsilon+1}\pi_0\biggl(\sum_{|A|-\text{odd}}q^A e_A\biggr)+\frac{1}{\epsilon-1}\pi_n\biggl(\sum_{|A|-\text{odd}}q^A e_A\biggr)\biggr)$$
в случае четного $n$ и
$$X=\frac{1}{2^{n-1}\epsilon}\biggl(\sum_{|A|-\text{odd}}q^A e_A-\frac{1}{\epsilon+1}\pi_0
\biggl(\sum_{|A|-\text{odd}}q^A e_A\biggr)-\frac{1}{\epsilon+1}\pi_n
\biggl(\sum_{|A|-\text{odd}}q^A e_A\biggr)\biggr)$$
в случае нечетного $n$.
\end{theorem}

\smallskip

\begin{proof} Имеем
$$\sum_{|A|-\text{odd}}e^A X e_A+\epsilon X \sum_{|A|-\text{odd}}e^A e_A=\sum_{|A|-\text{odd}}q^A e_A.$$
Отсюда в случае четного $n$ имеем
$$2^{n-1}\bigl((\epsilon+1)\pi_0(X)+(\epsilon-1)\pi_n(X)\bigr)+\epsilon 2^{n-1}\sum_{k=1}^{n-1} \pi_k(X) =\sum_{|A|-\text{odd}}q^A e_A,$$
а в случае нечетного $n$ "---
$$2^{n-1}\bigl((\epsilon+1)\pi_0(X)+(\epsilon+1)\pi_n(X)\bigr)+\epsilon 2^{n-1} \sum_{k=1}^{n-1} \pi_k(X)=\sum_{|A|-\text{odd}}q^A e_A.$$
Рассматривая всевозможные случаи, получаем утверждение теоремы.
\end{proof}

\smallskip

\Section{Cвертки по кватернионным типам}
Несложно подсчитать размерности подпространств \eqref{shirokov:kv} кватернионных типов 
$d_m(n)=\dim{{C}\!\ell}_{\overline m}(p,q)=\sum_k C_n^{4k+m}$, $m=0, 1, 2, 3$:
$$d_0(n)=2^{n-2}+2^{\frac{n-2}{2}}\cos{\frac{\pi n}{4}},\qquad d_1(n)=2^{n-2}+2^{\frac{n-2}{2}}\sin{\frac{\pi n}{4}},$$
$$d_2(n)=2^{n-2}-2^{\frac{n-2}{2}}\cos{\frac{\pi n}{4}},\qquad d_3(n)=2^{n-2}-2^{\frac{n-2}{2}}\sin{\frac{\pi n}{4}}.$$

Будем рассматривать следующие {\it свертки по кватернионным типам} (нормировать на множители $d_m(n)$ в дальнейшем иногда не будем):
\begin{gather}
F_{\overline 0}(U)=\frac{1}{d_0(n)}\sum_{A\in{\rm I}_ {\overline 0}} e_A U e^A,\qquad F_{\overline 1}(U)=\frac{1}{d_1(n)}\sum_{A\in{\rm I}_ {\overline 1}} e_A U e^A,\label{shirokov:svkvat}\\
F_{\overline 2}(U)=\frac{1}{d_2(n)}\sum_{A\in{\rm I}_ {\overline 2}} e_A U e^A,\qquad F_{\overline 3}(U)=\frac{1}{d_3(n)}\sum_{A\in{\rm I}_ {\overline 3}} e_A U e^A.\nonumber
\end{gather}

\smallskip

\hypertarget{shirokov:theoremSvKvat}{}
%\label{shirokov:theoremSvKvat} 
\begin{theorem}[8]Пусть $U_k$ "--- элемент алгебры Клиффорда ${{C}\!\ell}(p,q)$ ранга $k.$ 
При $k=1, \ldots, n-1$ свертки равны следующим величинам:
\begin{equation}
\begin{array}{l}
\displaystyle
\sum_{A\in{\rm I}_ {\overline 0}} e_A U_k e^A=2^{\frac{n-2}{2}}\cos\Bigl(\frac{\pi k}{2}-\frac{\pi n}{4}\Bigr)U_k,\\[2mm]
\displaystyle
\sum_{A\in{\rm I}_ {\overline 1}} e_A U_k e^A=(-1)^{k+1}2^{\frac{n-2}{2}}\sin\Bigl(\frac{\pi k}{2}-\frac{\pi n}{4}\Bigr)U_k,\\[2mm]
\displaystyle
\sum_{A\in{\rm I}_ {\overline 2}} e_A U_k e^A=-2^{\frac{n-2}{2}}\cos\Bigl(\frac{\pi k}{2}-\frac{\pi n}{4}\Bigr)U_k,\\[2mm]
\displaystyle
\sum_{A\in{\rm I}_ {\overline 3}} e_A U_k e^A=(-1)^{k}2^{\frac{n-2}{2}}\sin\Bigl(\frac{\pi k}{2}-\frac{\pi n}{4}\Bigr)U_k.
\end{array}
\label{shirokov:svkv}
\end{equation}

При $k=0$ и $k=n$ имеем $($для $m=0, 1, 2, 3)$
\begin{equation} 
\sum_{A\in{\rm I}_ {\overline m}} e_A  e^A=d_m(n)e,\qquad \sum_{A\in{\rm I}_ {\overline m}} e_A e^{1\ldots n} e^A=(-1)^{m(n+1)} d_m(n)e^{1\ldots n}.
\label{shirokov:svkv2}
\end{equation}
\end{theorem}

\smallskip

Заметим, что формулы \eqref{shirokov:svkv2} не являются частным случаем формул \eqref{shirokov:svkv}. 
Как следствие теоремы, имеем следующие формулы для сверток по кватернионным типам от произвольного элемента алгебры Клиффорда $U\in{{C}\!\ell}(p,q)$:
$$
\sum_{A\in{\rm I}_ {\overline 0}} 
e_A U e^A=\sum_{k=0}^3 2^{\frac{n-2}{2}}\cos\Bigl(\frac{\pi k}{2}-\frac{\pi n}{4}\Bigr)\pi_{\overline k}(U)+2^{n-2}\bigl(\pi_0(U)+\pi_n(U)\bigr),$$
\begin{multline}
\sum_{A\in{\rm I}_ {\overline 1}} e_A U e^A=\sum_{k=0}^3(-1)^{k+1}2^{\frac{n-2}{2}}\sin
\Bigl(\frac{\pi k}{2}-\frac{\pi n}{4}\Bigr)\pi_{\overline k}(U)+ \\ + 2^{n-2}\bigl(\pi_0(U)+(-1)^{n+1}\pi_n(U)\bigr),
\label{shirokov:sf}
\end{multline}
$$
\sum_{A\in{\rm I}_ {\overline 2}} e_A U e^A=\sum_{k=0}^3-2^{\frac{n-2}{2}}\cos\Bigl(\frac{\pi k}{2}-\frac{\pi n}{4}\Bigr)\pi_{\overline k}(U)+2^{n-2}\bigl(\pi_0(U)+\pi_n(U)\bigr),
$$
\begin{multline}
\sum_{A\in{\rm I}_ {\overline 3}} e_A U e^A=\sum_{k=0}^3(-1)^{k}2^{\frac{n-2}{2}}\sin\Bigl(\frac{\pi k}{2}-\frac{\pi n}{4}\Bigr)\pi_{\overline k}(U) +\\ +2^{n-2}\bigl(\pi_0(U)+(-1)^{n+1}\pi_n(U)\bigr).
\end{multline}

\smallskip

\begin{proof} Формулы для рангов $k=0$ и $k=n$ очевидны и следуют из формул для размерности соответствующих подпространств.
Для остальных $k=1, \ldots, n-1$, используя \eqref{shirokov:svert}, имеем
\begin{multline}
\sum_{A\in{\rm I}_ {\overline 0}} e_A U_k e^A=\biggl(\biggl(\sum_{i=0\mod 4}C_k^i\biggr)\biggl(\sum_{i=0\mod 4}C_{n-k}^i\biggr)- \\ 
- \biggl(\sum_{i=1\mod 4}C_k^i\biggr)\biggl(\sum_{i=3\mod 4}C_{n-k}^i\biggr)+
\biggl(\sum_{i=2\mod 4}C_k^i\biggr)\biggl(\sum_{i=2\mod 4}C_{n-k}^i\biggr)- \\ -
\biggl(\sum_{i=3\mod 4}C_k^i\biggr)\biggl(\sum_{i=1\mod 4}C_{n-k}^i\biggr)\biggr)U_k=
\\
=\bigl(d_0(k) d_{0}(n-k)-d_{1}(k) d_{3}(n-k)+d_{2}(k) d_{2}(n-k)-d_{3}(k) d_{1}(n-k)\bigr)U_k.
\end{multline}
Далее, пользуясь формулами для коэффициентов и тригонометрическими тождествами, получаем первое из утверждений теоремы.
Остальные случаи рассматриваются аналогично.
\end{proof}

\smallskip

Согласно доказанной теореме, свертки по мультииндексам фиксированных кватернионных типов сводятся к операциям проецирования на подпространства кватернионных типов и операциям проецирования на ранги $0$ и $n$. Однако в случае четного $n$ верно и обратное: можно выразить операции проецирования на подпространства кватернионных типов через четыре рассматриваемые свертки и две дополнительные свертки по рангам $0$ и $n$. А~именно, имеет место следующая теорема.

\smallskip

\begin{theorem}[9]Для произвольного элемента $U\in{{C}\!\ell}(p,q)$ в случае четного $n=p+q$ имеем
\begin{gather}
\pi_{\overline 0}(U)=2^{\frac{-n-2}{2}}\sum_{k=0}^3(-1)^k \cos\Bigl(\frac{\pi k}{2}-\frac{\pi n}{4}\Bigr)\sum_{A\in{\rm I}_ {\overline k}} e_A U e^A+2^{-2}(U+e_{1\ldots n}Ue^{1\ldots n}),
\\
\pi_{\overline 1}(U)=2^{\frac{-n-2}{2}}\sum_{k=0}^3\sin\Bigl(\frac{\pi k}{2}-\frac{\pi n}{4}\Bigr)\sum_{A\in{\rm I}_ {\overline k}} e_A U e^A+2^{-2}(U-e_{1\ldots n}Ue^{1\ldots n}),
\\
\pi_{\overline 2}(U)= 2^{\frac{-n-2}{2}}\sum_{k=0}^3(-1)^k \cos\Bigl(\frac{\pi k}{2}-\frac{\pi n}{4}\Bigr)\sum_{A\in{\rm I}_ {\overline k}} e_A U e^A+2^{-2}(U+e_{1\ldots n}Ue^{1\ldots n}),
\\
\pi_{\overline 3}(U)=2^{\frac{-n-2}{2}}\sum_{k=0}^3(-1)\sin\Bigl(\frac{\pi k}{2}-\frac{\pi n}{4}\Bigr)\sum_{A\in{\rm I}_ {\overline k}} e_A U e^A+2^{-2}(U-e_{1\ldots n}Ue^{1\ldots n}).
\end{gather}
\end{theorem}

\smallskip

\begin{proof}
Добавляем к рассматриваемым выше четырем уравнениям \eqref{shirokov:sf} два следующих уравнения:
$$eUe=\sum_{k=0}^3 \pi_{\overline k}(U),\qquad e_{1\ldots n}Ue^{1\ldots n}=\sum_{k=0}^3 (-1)^k \pi_{\overline k}(U).$$
Получающаяся квадратная матрица размера 6 рассматриваемой линейной системы уравнений имеет вид
{\small
\begin{eqnarray*}
\begin{pmatrix}
2^{\frac{n-2}{2}}\cos(\frac{\pi n}{4}) & 2^{\frac{n-2}{2}}\sin(\frac{\pi n}{4})& -2^{\frac{n-2}{2}}\cos(\frac{\pi n}{4}) & -2^{\frac{n-2}{2}}\sin(\frac{\pi n}{4})  & 2^{n-2} & 2^{n-2}\cr
2^{\frac{n-2}{2}}\sin(\frac{\pi n}{4}) & 2^{\frac{n-2}{2}}\cos(\frac{\pi n}{4})& -2^{\frac{n-2}{2}}\sin(\frac{\pi n}{4}) & -2^{\frac{n-2}{2}}\cos(\frac{\pi n}{4})  & 2^{n-2} & -2^{n-2}\cr
-2^{\frac{n-2}{2}}\cos(\frac{\pi n}{4}) & -2^{\frac{n-2}{2}}\sin(\frac{\pi n}{4})& 2^{\frac{n-2}{2}}\cos(\frac{\pi n}{4}) & 2^{\frac{n-2}{2}}\sin(\frac{\pi n}{4})  & 2^{n-2} & 2^{n-2}\cr
-2^{\frac{n-2}{2}}\sin(\frac{\pi n}{4}) & -2^{\frac{n-2}{2}}\cos(\frac{\pi n}{4})& 2^{\frac{n-2}{2}}\sin(\frac{\pi n}{4}) & 2^{\frac{n-2}{2}}\cos(\frac{\pi n}{4})  & 2^{n-2} & -2^{n-2}\cr
1 & 1& 1 & 1  & 0& 0\cr
1 & -1& 1 & -1  & 0 & 0\cr
\end{pmatrix}.\nonumber
\end{eqnarray*}
}

\noindent
Определитель этой матрицы равен $$2^{3n} \Bigl(\cos^2\Bigl(\frac{\pi n}{4}\Bigr)-\sin^2\Bigl(\frac{\pi n}{4}\Bigr)\Bigr).$$
В~случае нечетного $n$ определитель равен нулю и матрица необратима. 
В~случае $n=0\mod 4$ определитель равен $2^{3n}$, а в случае $n=2\mod 4$ определитель равен $-2^{3n}$ и матрица обратима. 
Можем записать в случае четного $n$, что определитель матрицы равен $(-1)^{{n}/{2}}2^{3n}$.

Обратная матрица имеет вид
{\small
\begin{eqnarray}
\begin{pmatrix}
2^{\frac{-n-2}{2}}\cos(\frac{\pi n}{4}) & -2^{\frac{-n-2}{2}}\sin(\frac{\pi n}{4})& -2^{\frac{-n-2}{2}}\cos(\frac{\pi n}{4}) & 2^{\frac{-n-2}{2}}\sin(\frac{\pi n}{4})  & 2^{-2} & 2^{-2}\cr
-2^{\frac{-n-2}{2}}\sin(\frac{\pi n}{4}) & 2^{\frac{-n-2}{2}}\cos(\frac{\pi n}{4})& 2^{\frac{-n-2}{2}}\sin(\frac{\pi n}{4}) & -2^{\frac{-n-2}{2}}\cos(\frac{\pi n}{4})  & 2^{-2} & -2^{-2}\cr
-2^{\frac{-n-2}{2}}\cos(\frac{\pi n}{4}) & 2^{\frac{-n-2}{2}}\sin(\frac{\pi n}{4})& 2^{\frac{-n-2}{2}}\cos(\frac{\pi n}{4}) & -2^{\frac{-n-2}{2}}\sin(\frac{\pi n}{4})  & 2^{-2} & 2^{-2}\cr
2^{\frac{-n-2}{2}}\sin(\frac{\pi n}{4}) & -2^{\frac{-n-2}{2}}\cos(\frac{\pi n}{4})& -2^{\frac{-n-2}{2}}\sin(\frac{\pi n}{4}) & 2^{\frac{-n-2}{2}}\cos(\frac{\pi n}{4})  & 2^{-2} & -2^{-2}\cr
2^{-n} & 2^{-n}& 2^{-n} & 2^{-n}  & 0& 0\cr
2^{-n} & -2^{-n}& 2^{-n} & -2^{-n}  & 0 & 0\cr
\end{pmatrix}.\nonumber
\end{eqnarray}

}

Таким образом, получаем в случае четного $n$ связь между проекциями на кватернионные типы и свертками, указанную в формулировке теоремы.
\end{proof}

\smallskip

Заметим, что в случае нечетного $n$ выписанная матрица необратима. 
Более того, имеем
 $$\sum_{A\in{\rm I}_ {\overline 0}} e_A U e^A=\sum_{A\in{\rm I}_ {\overline 3}} e_A U e^A,\qquad \sum_{A\in{\rm I}_ {\overline 2}} e_A U e^A=\sum_{A\in{\rm I}_ {\overline 1}} e_A U e^A,$$
 а значит
 \begin{multline}
e_A U e^A=2 \biggl(\sum_{A\in{\rm I}_ {\overline 0}} e_A U e^A+\sum_{A\in{\rm I}_ {\overline 2}} e_A U e^A\biggr)=2\biggl(\sum_{A\in{\rm I}_ {\overline 1}} e_A U e^A+\sum_{A\in{\rm I}_ {\overline 3}} e_A U e^A\biggr)=
\\
=2\biggl(\sum_{A\in{\rm I}_ {\overline 0}} e_A U e^A+\sum_{A\in{\rm I}_ {\overline 1}} e_A U e^A\biggr)=
2\biggl(\sum_{A\in{\rm I}_ {\overline 2}} e_A U e^A+\sum_{A\in{\rm I}_ {\overline 3}} e_A U e^A\biggr).
\end{multline}

Теперь рассмотрим системы коммутаторных уравнений по кватернионным типам, как это делалось выше для рангов, четных и нечетных сверток:
$$e^A X+\epsilon X e^A=q^A,\qquad |A|=m\mod 4,\qquad m=0, 1, 2, 3,\qquad \epsilon\in\C\setminus\{0\}.$$
В~силу громоздкости получающегося утверждения в общем случае рассмотрим в следующий теореме только случай $\epsilon=-1$ (случай коммутатора). 
В~других случаях система уравнений решается аналогично.

\smallskip

\begin{theorem}[10]Пусть элемент $X\in {{C}\!\ell}(p,q)$ алгебры Клиффорда размерности $p+q=n\geq 5$ удовлетворяет следующей системе уравнений для некоторых элементов $q^A\in{{C}\!\ell}(p,q){:}$
\begin{equation}
[e^A, X]=q^A,\quad \forall |A|=m\mod 4,\qquad m=0, 1, 2, 3.
\end{equation}
Тогда система либо не имеет ре\-ше\-ния$,$ либо имеет решение вида $($единст\-вен\-ное с точностью до прибавления указанных произвольных элементов фиксированных рангов$)$
$$
X=\frac{1}{2^{\frac{n-2}{2}}}\sum_{k=1}^{n-1}\frac{\pi_k\bigl(\sum_{|A|=0\mod 4}q^A e_A\bigr)}{\cos(\frac{\pi k}{2}- \frac{\pi n}{4})-\cos(\frac{\pi n}{4})-2^{\frac{n-2}{2}}}+U_0+U_n
$$
в случае $m=0,$
\begin{multline}
X=\frac{1}{2^{\frac{n-2}{2}}}\sum_{k=1}^{n-1}\frac{\pi_k\bigl(\sum_{|A|=1\mod 4}q^A e_A\bigr)}{(-1)^{k+1}\sin(\frac{\pi k}{2}-\frac{\pi n}{4})-\sin(\frac{\pi n}{4})-2^{\frac{n-2}{2}}}-
\\
-
\frac{\pi_n\bigl(\sum_{|A|=1\mod 4}q^A e_A\bigr)}{2^{n-1}+2^{\frac{n}{2}}\sin(\frac{\pi n}{4})}+U_0
\end{multline}
в случае $m=1$ и четного $n,$
$$X=\frac{1}{2^{\frac{n-2}{2}}}\sum_{k=1}^{n-1}\frac{\pi_k\bigl(\sum_{|A|=1\mod 4}q^A e_A\bigr)}{(-1)^{k+1}\sin(\frac{\pi k}{2}-\frac{\pi n}{4})-\sin(\frac{\pi n}{4})-2^{\frac{n-2}{2}}}+U_0+U_n$$
в случае $m=1$ и нечетного $n,$
$$X=\frac{1}{2^{\frac{n-2}{2}}}\sum_{k=1}^{n-1}\frac{\pi_k\bigl(\sum_{|A|=2\mod 4}q^A e_A\bigr)}{-\cos(\frac{\pi k}{2}-\frac{\pi n}{4})+\cos(\frac{\pi n}{4})-2^{\frac{n-2}{2}}}+U_0+U_n$$
в случае $m=2,$
\begin{multline}
X=\frac{1}{2^{\frac{n-2}{2}}}\sum_{k=1}^{n-1}\frac{\pi_k\bigl(\sum_{|A|=3\mod 4}q^A e_A\bigr)}{(-1)^{k}\sin(\frac{\pi k}{2}-\frac{\pi n}{4})+\sin(\frac{\pi n}{4})-2^{\frac{n-2}{2}}}- \\
- \frac{\pi_n\bigl(\sum_{|A|=3\mod 4}q^A e_A\bigr)}{2^{n-1}-2^{\frac{n}{2}}\sin(\frac{\pi n}{4})}+U_0
\end{multline}
в случае $m=3$ и четного $n,$
$$X=\frac{1}{2^{\frac{n-2}{2}}}\sum_{k=1}^{n-1}\frac{\pi_k\bigl(\sum_{|A|=3\mod 4}q^A e_A\bigr)}{(-1)^{k}\sin(\frac{\pi k}{2}-\frac{\pi n}{4})+\sin(\frac{\pi n}{4})-2^{\frac{n-2}{2}}}+U_0+U_n$$
в случае $m=3$ и нечетного $n.$
\end{theorem}

\smallskip

\begin{proof} Имеем для $m=0, 1, 2, 3$
$$\sum_{|A|=m\mod4}e^A X e_A-X \sum_{|A|=m\mod 4}e^A e_A=\sum_{|A|=m\mod 4}q^A e_A.$$

Пользуясь (\ref{shirokov:svkv}) и (\ref{shirokov:svkv2}), получаем для случая $m=0$
\begin{multline}
2^{\frac{n}{2}-1}\sum_{k=1}^{n-1}\cos\Bigl(\frac{\pi k}{2}-\frac{\pi n}{4}\Bigr)\pi_k(X)+
\Bigl(2^{n-2}+2^{\frac{n-2}{2}}\cos\Bigl(\frac{\pi n}{4}\Bigr)\Bigr)\pi_0(X) +
\\
+\Bigl(2^{n-2}+2^{\frac{n-2}{2}}\cos\Bigl(\frac{\pi n}{4}\Bigr)\Bigr)\pi_n(X)-
\Bigl(2^{n-2}+2^{\frac{n-2}{2}}\cos\Bigl(\frac{\pi n}{4}\Bigr)\Bigr)X=\sum_{|A|=0\mod 4}q^A e_A.
\end{multline}
%$$2^{\frac{n}{2}-1}\sum_{k=1}^{n-1}(-1)^{k+1}\sin(\frac{\pi k}{2}-\frac{\pi n}{4})\pi_k(X)+(2^{n-2}+2^{\frac{n-2}{2}}\sin(\frac{\pi n}{4}))\pi_0(X)+$$
%$$+(-1)^{n+1}(2^{n-2}+2^{\frac{n-2}{2}}\sin(\frac{\pi n}{4}))\pi_n(X)-(2^{n-2}+2^{\frac{n-2}{2}}\sin(\frac{\pi n}{4}))X=\sum_{|A|=1\mod 4}q^A e_A,$$
%$$-2^{\frac{n}{2}-1}\sum_{k=1}^{n-1}\cos(\frac{\pi k}{2}-\frac{\pi n}{4})\pi_k(X)+(2^{n-2}-2^{\frac{n-2}{2}}\cos(\frac{\pi n}{4}))\pi_0(X) +$$ $$+(2^{n-2}-2^{\frac{n-2}{2}}\cos(\frac{\pi n}{4}))\pi_n(X)-(2^{n-2}-2^{\frac{n-2}{2}}\cos(\frac{\pi n}{4}))X=\sum_{|A|=2\mod 4}q^A e_A,$$
%$$2^{\frac{n}{2}-1}\sum_{k=1}^{n-1}(-1)^{k}\sin(\frac{\pi k}{2}-\frac{\pi n}{4})\pi_k(X)+(2^{n-2}-2^{\frac{n-2}{2}}\sin(\frac{\pi n}{4}))\pi_0(X)+$$
%$$+(-1)^{n+1}(2^{n-2}-2^{\frac{n-2}{2}}\sin(\frac{\pi n}{4}))\pi_n(X)-(2^{n-2}-2^{\frac{n-2}{2}}\sin(\frac{\pi n}{4}))X=\sum_{|A|=3\mod 4}q^A e_A.$$
В случаях $m=1, 2, 3$ действуем аналогично.
Далее рассматриваем всевозможные случаи. Заметим, что выражения, стоящие в формулировке теоремы в знаменателях, не обращаются в ноль ни для какого $k$ при $n\geq 5$.
\end{proof}

\smallskip


Отметим, что можно сформулировать аналог этой теоремы и для случая $n\leq 4$. 
В~этом случае некоторые коэффициенты из доказательства предыдущей теоремы будут обнуляться для фиксированных рангов $k$. 
В таком случае решение будет содержать в качестве слагаемого вместо проекции на соответствующий ранг "--- произвольный элемент указанного ранга. Например, для $m=0$ коэффициент обнуляется при $n=4$ и $k=2$: 
$$\cos\Bigl(\frac{\pi k}{2}-\frac{\pi n}{4}\Bigr)-\cos\Bigl(\frac{\pi n}{4}\Bigr)-2^{\frac{n-2}{2}}=0,
$$ 
а значит, вместо проекции $\pi_2$ в качестве слагаемого будет присутствовать произвольный элемент $U_2$.

Заметим, что в случае размерностей $n\leq 3$ понятие кватернионного типа совпадает с понятием ранга элемента алгебры Клиффорда, поэтому все сводится к рассмотрению коммутаторных уравнений по рангам (см. рассуждения в~конце параграфа~\hyperlink{shir:par2}{2}).

\smallskip

\Section[n]{Заключение}
В настоящей статье представлены и доказаны утверждения для сверток в алгебрах Клиффорда разного вида, построенных с~помощью фиксированного базиса алгебры Клиффорда. 
Доказанные утверждения могут непосредственно применяться при решении уравнений в формализме алгебр Клиффорда, как это было сделано в работе \cite{shir:YM} при изучении уравнений Янга"--~Миллса и простейшего полевого уравнения.

Заметим, что во всех теоремах настоящей статьи мы могли рассматривать не генераторы алгебры Клиффорда $e^a$, а~произвольный набор элементов алгебры Клиффорда $\gamma^a\in{{C}\!\ell}(p,q)$, который удовлетворяет определяющим соотношениям $\gamma^a \gamma^b+ \gamma^b \gamma^a=2\eta^{ab}e$. 
Этот набор может порождать другой базис $\gamma^A$ алгебры Клиффорда ${{C}\!\ell}(p,q)$, 
но в~некоторых случаях нечетной размерности $n$ этот набор может не порождать новый базис ${{C}\!\ell}(p,q)$ (см. более подробно \cite{shir:gencontr}). 
В~работе \cite{shir:gencontr} рассматриваются обобщенные свертки, построенные по двум таким наборам $\gamma^a$, $\beta^a$.



\thanks{\textbf{Благодарности.} 
Работа выполнена при поддержке Российского научного фонда (РНФ 14--11--00687, Математический институт им.~В.~А.~Стеклова).
}

\thanks{{\bf ORCID}

{\bf Дмитрий Сергеевич Широков:} \url{http://orcid.org/0000-0002-2407-9601}


}


\RusRef

\begin{thebibliography}{99}

\Bibitem{shir:Clifford}
\paper Application of Grassmann's Extensive Algebra
\by    Clifford~W.~K.
\jour American Journal of Mathematics
\vol 1
\issue 4
\pages 350--358
\yr 1878
\crossref{10.2307/2369379}

\Bibitem{shir:Hamilton}
\paper II. On quaternions, or on a new system of imaginaries in algebra
\by    Hamilton~W.~R.
\jour Philosophical Magazine Series~3
\pages 489–495
\vol 25
\issue 163
\yr 1844
\crossref{10.1080/14786444408644923}

\Bibitem{shir:Grassmann}
\book  Die Lineale Ausdehnungslehre ein neuer Zweig der Mathematik 
\by    Grassmann~H.
\yr 1844
\totalpages xxxii+282
\publaddr Leipzig
\publ Verlag von Otto Wigand
\elink  Internet Archive Identifier: \href{https://archive.org/details/dielinealeausde00grasgoog}{dielinealeausde00grasgoog} 
\moreref
\book  Die Lineale Ausdehnungslehre ein neuer Zweig der Mathematik 
\by    Grassmann~H.
\yr 2012
\crossref{10.1017/CBO9781139237352}
\publ  Cambridge University Press
\publaddr Cambridge
\totalpages xxxii+282



\Bibitem{shir:Lipsch}
\book Untersuchungen \"uber die Summen von Quadraten
\by    Lipschitz~R.
\publ Max Cohen und Sohn
\publaddr Bonn
\totalpages 147
\yr 1886

\Bibitem{shir:Chevall}
\book Collected works
\bookvol 2
\volinfo The algebraic theory of spinors and Clifford algebras
\eds Pierre Cartier and Catherine Chevalley 
\publaddr Berlin
\publ Springer
\totalpages xiv+ 214 
\yr 1997
\by    Chevalley~C.
\zmath{https://zbmath.org/?q=an:0899.01032}


\Bibitem{shir:Dirac}
\paper The Quantum Theory of the Electron
\by    Dirac~P.~A.~M.
\jour Proc. R. Soc. (A)
\vol  117
\issue  778
\pages  610--624
\yr 1928
\crossref{10.1098/rspa.1928.0023}
\moreref
\crossref{10.1016/b978-0-08-006995-1.50017-x}
\paper The Quantum Theory of the Electron
\by    Dirac~P.~A.~M.
\inbook Special Theory of Relativity
\serial The Commonwealth and International Library: Selected Readings in Physics
\yr 1970
\pages 237–256


\Bibitem{shir:Hestenes}
\book  Clifford Algebra to Geometric Calculus. A Unified Language for Mathematics and Physics 
\by    Hestenes~D., Sobczyk~G.
\publaddr Reidel Publishing Company
\yr 1984
\totalpages 314


\RBibitem{shir:Marchuk}
\book Уравнения теории поля и алгебры Клиффорда
\by    Марчук~Н.~Г.
\publaddr Ижевск
\publ РХД
\yr 2009
\totalpages 304


\Bibitem{shir:tr1}
\paper   Computing Irreducible Representations of Groups 
\by Dixon~J.~D.
\jour  Math. Comp.
\yr 1970
\vol 24
\issue 111
\pages 707--712 
\crossref{10.2307/2004848}

\Bibitem{shir:tr2}
\paper Approximate representation theory of finite groups
\by    Babai~L.,  Friedl~K.
\jour Foundations of Computer Science
\yr 1991
\pages 733--742 
\crossref{10.1109/sfcs.1991.185442}

\Bibitem{shir:aacacontr}
\by Shirokov~D.~S.
\preprint Method of averaging in Clifford algebras
\yr 2015 
\arxiv \href{http://arxiv.org/abs/1412.0246}{1412.0246} [math-ph]
\totalpages 15



\Bibitem{shir:YM}
\preprint New class of gauge invariant solutions of Yang--Mills equations
\by Marchuk~N.~G., Shirokov~D.~S.
\yr 2014 
\arxiv \href{http://arxiv.org/abs/1406.6665}{1406.6665} [math-ph]
\totalpages 35


\Bibitem{shir:gencontr}
\by Shirokov~D.~S.
\preprint Method of generalized contractions and Pauli's theorem in Clifford algebras
\yr 2014 
\arxiv \href{http://arxiv.org/abs/1409.8163}{1409.8163} [math-ph]
\totalpages 14


\Bibitem{shir:Pauli}
\paper Contributions math\'ematiques \`a la th\'eorie des matrices de Dirac
\by    Pauli~W.
\jour Annales de l'institut Henri Poincar\'e 
\yr 1936
\vol 6
\issue 2
\pages 109--136




\RBibitem{shir:Pauli1}
\by Широков~Д.~С.
\paper Обобщение теоремы Паули на случай алгебр Клиффорда
\jour Докл. РАН
\vol 440
\issue 5 
\yr 2011
\pages 1–4  
\zmath{https://zbmath.org/?q=an:1237.81095}
%\mathscinet{http://www.ams.org/mathscinet-getitem?mr=2839068}
%D. S. Shirokov, “Extension of Pauli's theorem to Clifford algebras”, Dokl. Math., 84:2 (2011), 699–701  crossref  zmath  isi  scopus


\RBibitem{shir:Pauli2}
\by Широков~Д.~С.
\paper Теорема Паули при описании $n$-мерных спиноров в формализме алгебр Клиффорда
\jour ТМФ
\yr 2013
\vol 175
\issue 1
\pages 11--34
\mathnet{http://mi.mathnet.ru/tmf8384}
\crossref{10.4213/tmf8384}
\mathscinet{http://www.ams.org/mathscinet-getitem?mr=3172130}
\zmath{https://zbmath.org/?q=an:1286.81070}
\adsnasa{http://adsabs.harvard.edu/cgi-bin/bib_query?2013TMP...175..454S}
%\transl
%\jour Theoret. and Math. Phys.
%\yr 2013
%\vol 175
%\issue 1
%\pages 454--474
%\crossref{http://dx.doi.org/10.1007/s11232-013-0038-9}
%\isi{http://gateway.isiknowledge.com/gateway/Gateway.cgi?GWVersion=2&SrcApp=PARTNER_APP&SrcAuth=LinksAMR&DestLinkType=FullRecord&DestApp=ALL_WOS&KeyUT=000318805200002}
%\scopus{http://www.scopus.com/record/display.url?origin=inward&eid=2-s2.0-84877350521}


\RBibitem{shir:Pauli3}
\by Широков~Д.~С.
\paper Использование обобщённой теоремы Паули для нечётных элементов алгебры Клиффорда 
для анализа связей между спинорными и ортогональными группами произвольных размерностей
\jour Вестн. Сам. гос. техн. ун-та. Сер. Физ.-мат. науки
\yr 2013
\issue 1(30)
\pages 279--287
\mathnet{http://mi.mathnet.ru/vsgtu1176}
\crossref{10.14498/vsgtu1176}


\Bibitem{shir:Marchuk:Shirokov}
\paper Unitary spaces on Clifford algebras
\by Marchuk~N.~G., Shirokov~D.~S.
\jour Adv. Appl. Clifford Algebras
\yr 2008
\vol 18
\issue 2
\pages 237--254
\crossref{10.1007/s00006-008-0066-y}
\arxiv \href{http://arxiv.org/abs/0705.1641}{0705.1641}  [math-ph]

\Bibitem{shir:Lounesto}
\by Lounesto~P.
\publ Cambridge University Press
\publaddr Cambridge
\serial 	London Mathematical Society Lecture Note Series	
\vol 239
\book Clifford Algebras and Spinors
\yr 1997
\totalpages ix+306 
\zmath{https://zbmath.org/?q=an:0887.15029}
\moreref
\by Lounesto~P.
\book Clifford Algebras and Spinors \rm (second edition)
\publ Cambridge University Press
\publaddr Cambridge 
\crossref{10.1017/cbo9780511526022}
\serial 	London Mathematical Society Lecture Note Series	
\vol 286
\yr 2001
\totalpages ix+338
\zmath{https://zbmath.org/?q=an:01591940}


\RBibitem{shir:quat1}
\by Широков~Д.~С.
\paper Классификация элементов алгебр Клиффорда по кватернионным типам
\jour ДАН
\yr 2009
\vol 427
\issue 6
\pages 758--760
\mathscinet{http://www.ams.org/mathscinet-getitem?mr=2590703}
\zmath{https://zbmath.org/?q=an:1183.15021}
\elib{http://elibrary.ru/item.asp?id=12803924}
%\transl
%\by Shirokov~D.\,S.
%\paper Classification of elements of clifford algebras according to quaternionic types
%\jour Doklady Mathematics
%\yr 2009
%\vol 80
%\issue 1
%\pages 610--612
%\crossref{10.1134/S1064562409040401}
%\mathscinet{http://www.ams.org/mathscinet-getitem?mr=2590703}
%\zmath{https://zbmath.org/?q=an:1183.15021}
%\isi{http://gateway.isiknowledge.com/gateway/Gateway.cgi?GWVersion=2&SrcApp=PARTNER_APP&SrcAuth=LinksAMR&DestLinkType=FullRecord&DestApp=ALL_WOS&KeyUT=000269308500040}
%\elib{http://elibrary.ru/item.asp?id=15299989}


\Bibitem{shir:quat2}
\paper Quaternion typification of Clifford algebra elements
\by Shirokov~D.~S.
\jour Adv. Appl. Clifford Algebras
\vol 22
\issue 1
\pages 243-256
\yr 2012
\crossref{10.1007/s00006-011-0288-2}

\Bibitem{shir:quat3}
\paper Development of the method of quaternion typification of clifford algebra elements
\by Shirokov~D.~S.
\jour Adv. Appl. Clifford Algebras
\vol 22
\issue 2
\pages 483--497
\yr 2012
\crossref{10.1007/s00006-011-0304-6}
\arxiv \href{http://arxiv.org/abs/0903.3494}{0903.3494}  [math-ph]

\end{thebibliography}

\RuDate


\smallskip

\makeentitle



\thanks{\textbf{Acknowledgments.} 
This work was supported by Russian Science Foundation (RSF 14--11--00687, Steklov Mathematical Institute).
}

\thanks{{\bf ORCID}

{\bf Dmitry S. Shirokov:} \url{http://orcid.org/0000-0002-2407-9601}


}


\EngRef

\begin{thebibliography}{99}

\Bibitem{shir:Clifford:eng}
\paper Application of Grassmann's Extensive Algebra
\by    Clifford~W.~K.
\jour American Journal of Mathematics
\vol 1
\issue 4
\pages 350--358
\yr 1878
\crossref{10.2307/2369379}

\Bibitem{shir:Hamilton:eng}
\paper II. On quaternions, or on a new system of imaginaries in algebra
\by    Hamilton~W.~R.
\jour Philosophical Magazine Series~3
\pages 489–495
\vol 25
\issue 163
\yr 1844
\crossref{10.1080/14786444408644923}

\Bibitem{shir:Grassmann:eng}
\book  Die Lineale Ausdehnungslehre ein neuer Zweig der Mathematik 
\by    Grassmann~H.
\yr 1844
\totalpages xxxii+282
\publaddr Leipzig
\publ Verlag von Otto Wigand
\elink  Internet Archive Identifier: \href{https://archive.org/details/dielinealeausde00grasgoog}{dielinealeausde00grasgoog} 
\moreref
\book  Die Lineale Ausdehnungslehre ein neuer Zweig der Mathematik 
\by    Grassmann~H.
\yr 2012
\crossref{10.1017/CBO9781139237352}
\publ  Cambridge University Press
\publaddr Cambridge
\totalpages xxxii+282



\Bibitem{shir:Lipsch:eng}
\book Untersuchungen \"uber die Summen von Quadraten
\by    Lipschitz~R.
\publ Max Cohen und Sohn
\publaddr Bonn
\totalpages 147
\yr 1886

\Bibitem{shir:Chevall:eng}
\book Collected works
\bookvol 2
\volinfo The algebraic theory of spinors and Clifford algebras
\eds Pierre Cartier and Catherine Chevalley 
\publaddr Berlin
\publ Springer
\totalpages xiv+ 214 
\yr 1997
\by    Chevalley~C.
\zmath{https://zbmath.org/?q=an:0899.01032}


\Bibitem{shir:Dirac:eng}
\paper The Quantum Theory of the Electron
\by    Dirac~P.~A.~M.
\jour Proc. R. Soc. (A)
\vol  117
\issue  778
\pages  610--624
\yr 1928
\crossref{10.1098/rspa.1928.0023}
\moreref
\crossref{10.1016/b978-0-08-006995-1.50017-x}
\paper The Quantum Theory of the Electron
\by    Dirac~P.~A.~M.
\inbook Special Theory of Relativity
\serial The Commonwealth and International Library: Selected Readings in Physics
\yr 1970
\pages 237–256


\Bibitem{shir:Hestenes:eng}
\book  Clifford Algebra to Geometric Calculus. A Unified Language for Mathematics and Physics 
\by    Hestenes~D., Sobczyk~G.
\publaddr Reidel Publishing Company
\yr 1984
\totalpages 314


\Bibitem{shir:Marchuk:eng}
\lang In Russian
\book Uravneniia teorii polia i algebry Klifforda \rm [Equations of Field Theory and Clifford Algebras]
\by    Marchuk~N.~G.
\publaddr Izhevsk
\publ Regulyarnaya i Khaoticheskaya Dinamika
\yr 2009
\totalpages 304

\Bibitem{shir:tr1:eng}
\paper   Computing Irreducible Representations of Groups 
\by Dixon~J.~D.
\jour  Math. Comp.
\yr 1970
\vol 24
\issue 111
\pages 707--712 
\crossref{10.2307/2004848}

\Bibitem{shir:tr2:eng}
\paper Approximate representation theory of finite groups
\by    Babai~L.,  Friedl~K.
\jour Foundations of Computer Science
\yr 1991
\pages 733--742 
\crossref{10.1109/sfcs.1991.185442}

\Bibitem{shir:aacacontr:eng}
\by Shirokov~D.~S.
\preprint Method of averaging in Clifford algebras
\yr 2015 
\arxiv \href{http://arxiv.org/abs/1412.0246}{1412.0246} [math-ph]
\totalpages 15


\Bibitem{shir:YM:eng}
\preprint New class of gauge invariant solutions of Yang--Mills equations
\by Marchuk~N.~G., Shirokov~D.~S.
\yr 2014 
\arxiv \href{http://arxiv.org/abs/1406.6665}{1406.6665} [math-ph]
\totalpages 35


\Bibitem{shir:gencontr:eng}
\by Shirokov~D.~S.
\preprint Method of generalized contractions and Pauli's theorem in Clifford algebras
\yr 2014 
\arxiv \href{http://arxiv.org/abs/1409.8163}{1409.8163} [math-ph]
\totalpages 14


\Bibitem{shir:Pauli:eng}
\paper Contributions math\'ematiques \`a la th\'eorie des matrices de Dirac
\by    Pauli~W.
\jour Annales de l'institut Henri Poincar\'e 
\yr 1936
\vol 6
\issue 2
\pages 109--136




\RBibitem{shir:Pauli1:eng}
\zmath{https://zbmath.org/?q=an:1237.81095}
%\mathscinet{http://www.ams.org/mathscinet-getitem?mr=2839068}
\by Shirokov~D.~S.
\paper Extension of Pauli's theorem to Clifford algebras
\jour  Dokl. Math.
\vol 84
\issue 2 
\yr 2011
\pages 699–701  
\crossref{10.1134/s1064562411060329}



\Bibitem{shir:Pauli2:eng}
\paper Pauli theorem in the description of n -dimensional spinors in the Clifford algebra formalism
\by Shirokov~D.~S.
\jour Theoret. and Math. Phys.
\yr 2013
\vol 175
\issue 1
\pages 454--474
\crossref{10.1007/s11232-013-0038-9}
\isi{http://gateway.isiknowledge.com/gateway/Gateway.cgi?GWVersion=2&SrcApp=PARTNER_APP&SrcAuth=LinksAMR&DestLinkType=FullRecord&DestApp=ALL_WOS&KeyUT=000318805200002}
\scopus{http://www.scopus.com/record/display.url?origin=inward&eid=2-s2.0-84877350521}


\Bibitem{shir:Pauli3:eng}
\lang In Russian
\by Shirokov~D.~S.
\paper The use of the generalized Pauli's theorem for odd elements of Clifford algebra
to analyze relations between spin and orthogonal groups of arbitrary dimensions
\jour Vestn. Samar. Gos. Tekhn. Univ. Ser. Fiz.-Mat. Nauki \rm [J. Samara State Tech. Univ., Ser. Phys. \& Math. Sci.]
\yr 2013
\issue 1(30)
\pages 279--287
\crossref{10.14498/vsgtu1176}


\Bibitem{shir:Marchuk:Shirokov:eng}
\paper Unitary spaces on Clifford algebras
\by Marchuk~N.~G., Shirokov~D.~S.
\jour Adv. Appl. Clifford Algebras
\yr 2008
\vol 18
\issue 2
\pages 237--254
\crossref{10.1007/s00006-008-0066-y}
\arxiv \href{http://arxiv.org/abs/0705.1641}{0705.1641}  [math-ph]

\Bibitem{shir:Lounesto:eng}
\by Lounesto~P.
\publ Cambridge University Press
\publaddr Cambridge
\serial 	London Mathematical Society Lecture Note Series	
\vol 239
\book Clifford Algebras and Spinors
\yr 1997
\totalpages ix+306 
\zmath{https://zbmath.org/?q=an:0887.15029}
\moreref
\by Lounesto~P.
\book Clifford Algebras and Spinors \rm (second edition)
\publ Cambridge University Press
\publaddr Cambridge 
\crossref{10.1017/cbo9780511526022}
\serial 	London Mathematical Society Lecture Note Series	
\vol 286
\yr 2001
\totalpages ix+338
\zmath{https://zbmath.org/?q=an:01591940}


\Bibitem{shir:quat1:eng}
\by Shirokov~D.~S.
\paper Classification of elements of clifford algebras according to quaternionic types
\jour Doklady Mathematics
\yr 2009
\vol 80
\issue 1
\pages 610--612
\crossref{10.1134/S1064562409040401}
\mathscinet{http://www.ams.org/mathscinet-getitem?mr=2590703}
\zmath{https://zbmath.org/?q=an:1183.15021}
\isi{http://gateway.isiknowledge.com/gateway/Gateway.cgi?GWVersion=2&SrcApp=PARTNER_APP&SrcAuth=LinksAMR&DestLinkType=FullRecord&DestApp=ALL_WOS&KeyUT=000269308500040}
\elib{http://elibrary.ru/item.asp?id=15299989}


\Bibitem{shir:quat2:eng}
\paper Quaternion typification of Clifford algebra elements
\by Shirokov~D.~S.
\jour Adv. Appl. Clifford Algebras
\vol 22
\issue 1
\pages 243-256
\yr 2012
\crossref{10.1007/s00006-011-0288-2}

\Bibitem{shir:quat3:eng}
\paper Development of the method of quaternion typification of clifford algebra elements
\by Shirokov~D.~S.
\jour Adv. Appl. Clifford Algebras
\vol 22
\issue 2
\pages 483--497
\yr 2012
\crossref{10.1007/s00006-011-0304-6}
\arxiv \href{http://arxiv.org/abs/0903.3494}{0903.3494}  [math-ph]

\end{thebibliography}

\EngDate

\newpage
\Russian

%
%\end{document}
